\section{Related Work}
\label{sec:related}

\subsection{U-space and 4D Trajectory-Based Operations}

The SESAR Master Plan situates U-space within the broader transition toward trajectory-based operations (TBO) and digitally enabled air traffic management~\cite{sesar_master_plan_2025}. Comparable developments in FAA NextGen emphasize negotiated 4D trajectories to improve predictability and capacity~\cite{enea_porretta_4d_tbo_comparison}. These initiatives highlight time-constrained trajectory management and interoperable services as key enablers of scalable urban air traffic integration.

Conflict management in discretized 4D space--time representations has been demonstrated for multi-UAS operations~\cite{acevedo_4d_grids_u_space}. In parallel, U-space service architectures have been proposed that combine monitoring, traffic management, and in-flight service functions in modular system concepts~\cite{uspace_inflight_services_architecture}. These works provide important conceptual foundations for service-oriented airspace management, but comparatively few studies report computational performance measurements for deployable planning pipelines.

\subsection{UAS Path Planning and Conflict Detection}

Corridors are only one of several airspace models proposed for low-altitude urban operations. Alternative concepts include layered airspace structures, zone-based concepts, and largely unstructured free-flight airspace; the Metropolis project related these structuring options to achievable capacity under extreme traffic densities~\cite{sunil2015metropolis,schneider2014metropolis}, comparative performance assessments of different airspace management concepts for UAM have been reported more recently~\cite{abdellaoui2025_airspace_uam}, and DLR's UAM air-traffic-management research provides a broader review of these developments~\cite{schuchardt2023_atm_uam_review}. A corridor-based model is adopted in this work because it maps naturally onto a graph representation suitable for service-oriented route planning, offers predictable and procedurally simple route structures for early UAM implementation, and is envisioned by several operators for initial corridor-based operations within U-space; the trade-off is reduced routing flexibility and potentially lower capacity compared with less structured airspace concepts.

Search-based planning methods remain widely used for UAS routing in structured airspace. In corridor-based airspace models, the routing problem can naturally be represented as a graph, where nodes correspond to vertiports or route intersections and edges represent available corridor connections. Classical shortest-path methods and heuristic search algorithms such as A* are therefore suitable choices for computing lateral routes in such networks.

A* is particularly attractive in corridor networks because an admissible distance-based heuristic can guide the search toward the destination while preserving shortest-path optimality~\cite{sahawneh2016_3d_path_planning,a_star_2d3d_comparison_utm}. However, graph search alone does not produce a complete 4D flight trajectory. The selected lateral route must still be converted into a time-parameterized trajectory that respects flight-phase structure, altitude transitions, and aircraft-performance assumptions.

Conflict detection and resolution in multi-dimensional trajectory spaces have been extensively studied within trajectory-based operations frameworks~\cite{NDMap,enea_porretta_4d_tbo_comparison}. These approaches formalize separation assessment in time-dependent state spaces and provide a foundation for integrated 4D conflict management. Nevertheless, routing, trajectory generation, and deconfliction are often evaluated as separate algorithmic components rather than as parts of a service-oriented planning pipeline.

\subsection{Microservice Architectures in UAS Applications}

Microservice architectures have gained prominence in distributed software engineering due to their modularity, separation of concerns, and potential for independent scalability~\cite{dragoni2017microservices}. Applications in UAS-related domains demonstrate benefits for maintainability, simulation integration, and flexible deployment~\cite{matlekovic2020uav}.

For U-space planning systems, such an architectural style is attractive because topology provision, route planning, trajectory generation, deconfliction, and authorization can be assigned to separate service responsibilities. However, the quantitative evaluation of such integrated planning stacks remains limited. In particular, there is still a need to understand which stage of the planner pipeline dominates latency when graph-based routing is combined with 4D trajectory generation.

\medskip
\noindent
\textbf{Research Gap.}
Existing literature establishes theoretical and architectural foundations for U-space integration, UAS routing, and trajectory-based conflict management. However, few works report service-oriented latency decomposition for a deployable 4D path-planning pipeline or evaluate conflict checking and re-planning within the same service-oriented prototype. In particular, runtime breakdowns separating graph-based route search, waypoint synthesis, and kinematically feasible trajectory generation are rarely quantified. This paper addresses that gap by evaluating planner-stage computational performance and a preliminary NDMap-based lateral re-routing loop on an exemplary vertiport corridor network.
